
\documentclass[12pt]{article}
\usepackage{geometry}
\geometry{margin=1in}
\title{ECE 379K Design Document}
\date{}

\begin{document}
\maketitle

\section{Document Metadata}
Provide project title, team members, course info, date, and version. Add a revision history table and role descriptions (Technical Lead, Project Manager, Concurrency Specialist, etc.). Include schedules of what role each team member will play throughout the semester (if you are all going to share each role, state that here).

\section{Executive Summary}
Summarize the project in 1--2 paragraphs, describing the problem being solved, the overall approach, and the intended outcome.

\section{System Overview}
Include a polished architecture diagram and execution model. Explain major design decisions and trade-offs, including how tasks are split between CPU and GPU.

\section{Functional Requirements}
Provide a numbered list of functional requirements. Each should be specific, measurable, and testable.

\section{Non-Functional Requirements}
Detail performance, scalability, reliability, maintainability, and portability expectations. Provide measurable targets where possible.

\section{Detailed Design}
Expand on the module decomposition. For each module, describe its responsibilities, inputs/outputs, and interactions. Include diagrams, pseudocode, and data structures where appropriate.

\section{Performance \& Concurrency Plan}
Explain how work is partitioned across CPU and GPU. Describe concurrency strategies, bottleneck mitigation, load balancing, and plans for benchmarking.

\section{Testing \& Validation}
Describe your strategy for unit, integration, performance, and regression testing. Identify specific test cases and how success will be measured.

\section{Risks \& Mitigations}
Provide a risk table listing risks, their likelihood, impact, and mitigation strategies. Include contingency plans for the most critical risks.

\section{Project Management}
Detail the project timeline, milestones with dates, and responsibilities. Describe workflow practices such as Git branching strategy, code reviews, and CI/CD pipeline.

\section{References}
List all relevant references, including research papers, documentation, and libraries used. Follow a consistent citation format.

\section{Appendix}
Include supporting materials such as diagrams, pseudocode, glossary of terms, and any other resources that support your design.

\end{document}
